\section{Evidência de IMP-NOTIF-004 (expiração determinada na emissão)}

\textbf{Feature:} \texttt{IMP-NOTIF-004} --- expiração e conjunto ativo.

\textbf{Estado anterior:} \texttt{DIVERGENTE}. \texttt{adicionarNotificacaoTx}
aplicava 30 dias a \emph{todos} os tipos, contrariando RN-M13-03 (``não se
atribui duração arbitrária aos demais tipos'' e \texttt{ESCASSEZ\_ESTOQUE} não
expira). O enum do schema também divergia (\texttt{REAGENTE\_ESCASSO}; sem
\texttt{Bolsista}/\texttt{Chefe\_Geral}; sem \texttt{expira\_em}).

\textbf{Estado depois:} emissão, conjunto ativo e abertura revalidada do alvo
estão convergidos; a feature passa a \texttt{NÃO DIVERGENTE}.

\subsection{Contrato}

Fontes: RN-M13-03; CUE \texttt{\#M13Tipo} (inclui \texttt{ESCASSEZ\_ESTOQUE}),
\texttt{\#M13PapelDestinatario} e \texttt{\#M13Expiracao}
(\texttt{expira\_em = null} $\Rightarrow$ ativo; \texttt{ESCASSEZ\_ESTOQUE}
$\Rightarrow$ null).

Propriedades: sem prazo determinado na fonte, \texttt{expira\_em = null};
prazo explícito da fonte é preservado; \texttt{ESCASSEZ\_ESTOQUE} grava
\texttt{null} e rejeita prazo.

\subsection{Mudança}

\begin{itemize}
  \item \path{functions/src/notificacoes.ts}: \texttt{adicionarNotificacaoTx}
  deixou de impor 30 dias; grava \texttt{dados.expira\_em ?? null} e força
  \texttt{null} para \texttt{ESCASSEZ\_ESTOQUE} (rejeitando prazo com
  \texttt{invalid-argument});
  \item \path{functions/src/schemas/notificacoes.schema.ts}: novo
  \texttt{expira\_em}; \texttt{REAGENTE\_ESCASSO} $\rightarrow$
  \texttt{ESCASSEZ\_ESTOQUE}; adicionados \texttt{Bolsista} e
  \texttt{Chefe\_Geral} ao papel destinatário.
  \item \path{functions/src/notificacoes.ts}: \texttt{resolverDestinoNotificacao}
  relê autoridade M9, destinatário, expiração e ACL corrente do alvo antes de
  devolver rota; perda de vínculo ou alvo expirado retorna resposta neutra.
  \item \path{frontend/src/components/layout/NotificacoesDropdown.tsx} e
  \path{frontend/src/app/notificacoes/page.tsx}: abertura aguarda a validação
  server-side e o relógio do conjunto ativo é atualizado a cada 30 segundos.
\end{itemize}

\subsection{Testes}

\small
\begin{longtable}{@{}p{0.30\textwidth}p{0.52\textwidth}@{}}
\toprule
ID & Resultado \\
\midrule
\endhead
TEST-INT-NOTIF-M13-027 & Sem prazo de fonte grava \texttt{expira\_em=null} \\
TEST-INT-NOTIF-M13-028 & Prazo determinado da fonte é preservado \\
TEST-INT-NOTIF-M13-029 & \texttt{ESCASSEZ\_ESTOQUE} grava \texttt{expira\_em=null} \\
TEST-INT-NOTIF-M13-030 & \texttt{ESCASSEZ\_ESTOQUE} com prazo é rejeitada \\
TEST-INT-NOTIF-M13-031--033 & abertura revalida vínculo atual e mantém alvo expirado neutro \\
\bottomrule
\end{longtable}
\normalsize

RED: os quatro casos falharam contra a emissão anterior (30 dias universal).
GREEN sob Node 20.19.6: notificações 41/41 (integração 30/30 e Rules 11/11);
unit 71/71; Rules 77/77; integração completa 160 passam / 14 falham (legadas).

\textbf{Estado final:} \texttt{IMP-NOTIF-004} = \texttt{NÃO DIVERGENTE};
backend, caixa UI-12 e revalidação server-side passaram sob Node 20.
